\chapter{Introducción}

\section{Objetivo General}

Desarrollar un prototipo funcional que permita transformar datos operativos de tiendas de e-commerce en recomendaciones accionables para mejorar la toma de decisiones comerciales en la Ciudad Autónoma de Buenos Aires, Argentina, durante el año 2026.

\section{Objetivos Específicos}

El presente proyecto se plantea los siguientes objetivos específicos:

\begin{enumerate}
    \item Optimizar la eficiencia operativa de los comercios mediante la automatización del procesamiento de datos de ventas e inventario.

    \item Incrementar el volumen de ventas a través de la identificación de patrones de comportamiento que deriven en estrategias comerciales precisas.

    \item Mejorar el nivel de engagement y la experiencia del usuario final mediante la generación de recomendaciones personalizadas y oportunas.

    \item Validar la utilidad del sistema comparando la capacidad de respuesta de los usuarios ante recomendaciones directas versus dashboards descriptivos tradicionales.
\end{enumerate}

\section{Alcance}

El proyecto contempla el diseño y desarrollo de un prototipo funcional con cobertura geográfica limitada a la Ciudad Autónoma de Buenos Aires, Argentina.

\subsection{Componentes Incluidos}

\begin{enumerate}
    \item Módulo de gestión de usuarios (registro y autenticación/login).

    \item Adquisición y procesamiento de datos de e-commerce a través de integración con API de Tiendanube y carga manual de archivos CSV.

    \item Desarrollo de un motor de análisis para la detección de patrones y oportunidades comerciales, con enfoque en rubros específicos (gastronomía, indumentaria y electrónica).

    \item Construcción de una interfaz responsiva que presente recomendaciones accionables, optimizada para diversos dispositivos.

    \item Implementación de un mecanismo de retroalimentación para evaluar el impacto de las recomendaciones ejecutadas.
\end{enumerate}

\subsection{Componentes Excluidos}

El foco estará en demostrar la viabilidad de generar valor comercial, priorizando la claridad de las recomendaciones. Quedan fuera del alcance:

\begin{enumerate}
    \item Integraciones múltiples con plataformas externas más allá de Tiendanube.

    \item Despliegue productivo a escala industrial.

    \item Procesamiento de datos en tiempo real.
\end{enumerate}
